\documentclass[12pt]{article}
\usepackage[utf8]{inputenc}
\usepackage[T1]{fontenc}
\usepackage{hyperref}
\title{Buksa Language Resources for K-12 Curriculum Development}
\author{Compiled by Claude Sonnet 5 (claude-sonnet-5) for Andr\'as Kornai}
\date{}
\begin{document}
\maketitle

\section{Introduction}
Buksa (also Buksari, Bhoksa; ISO 639-3 tkb) is spoken mainly in Udham Singh Nagar and Nainital districts of Uttarakhand, India. Wikipedia gives about 59,000 speakers (2011 census), and the state tribal welfare site gives a community population of 46,771. Classification is disputed across sources: Indo-Aryan, Tharu-related, or a Hindi dialect. The search found very little dedicated material. This dossier is short because little exists in what I found.

\section{Primary Educational Resources}
No primers, textbooks or school programmes for Buksa turned up in the search. Wikipedia states the language has no written literature, only an oral tradition of folktales and folk songs.

\section{Linguistic and Reference Materials}
\begin{itemize}
\item Wikipedia, ``Buksa language'': \url{https://en.wikipedia.org/wiki/Buksa_language}. It cites Jagdish Pant's entry ``Buksa/Buksari'' in \emph{The Languages of Uttarakhand} (People's Linguistic Survey of India, vol.\ 30, 2015). I did not locate that volume online.
\item Wikipedia, ``Bhoksa people'': \url{https://en.wikipedia.org/wiki/Bhoksa_people}.
\end{itemize}
I found no dictionary or grammar.

\section{Religious and Community Resources}
\begin{itemize}
\item Uttarakhand Tribal Welfare Department, ``Buksa Tribes'': \url{https://tribalwelfare.uk.gov.in/page/buksa-tribes}. The Buksa are a Scheduled Tribe and a Particularly Vulnerable Tribal Group (see \url{https://tribalwelfare.uk.gov.in/page/pvgt-list}).
\end{itemize}

\section{Digital and Audio Resources}
No audio, video or corpus resources were found in the search.

\section{Curriculum Development Context}
There is no known written standard or school use. A curriculum would first need community-led documentation: recorded oral literature, an agreed orthography (Devanagari is the script listed), and a word list. Fieldwork by or with the Buksa community is the realistic starting point.
\end{document}